\section{Finite blocks and prediction}
\label{sec:blocks}

The signal from Lemma~\ref{lem:refresh-transmit} can be arbitrarily
weak. We use many such signals, holding one centering function fixed
throughout each block. For a block with $n\lambda^2$ comparable to
$\delta$, the square-root integral loses $O(\delta^2)$, while the
estimator has mean-square error $O(1/\delta)$. Averaging the estimators
from $K$ blocks gives mean-square error $O((K\delta)^{-1})$ with total
loss $O(K\delta^2)$. Thus $K\delta$ can grow while $K\delta^2$ tends
to zero.

\begin{proposition}[Prediction of a tagged packet sum]
\label{prop:prediction}
Let $S$ be an intermediate map, expressed in a fixed physical
volume-preserving coordinate chart.  Let
\[
 F_r=\sum_{m\in\mathcal M_r} A_{r,m}(u,a)\ee(m\theta),
 \qquad F=\sum_r F_r,
\]
be a finite list of colors of smooth packets supported compactly in the
good arc.  Suppose that the finite sets
$\mathcal M_r\subset\Z\setminus\{0\}$ are pairwise disjoint, that
$|A_{r,m}(u,a)|$ is independent of $a$, and that the list is invariant
under complex conjugation: $F_{r^*}=\overline{F_r}$.  In particular, $F$
is real.  Let $p_r$ be the corresponding spectral densities, put
$p=\sum_r p_r$, and suppose that
\[
 \int_\T p(z)\,\dd z\le \frac32.
\]
Assign constants $b_r\in\C$ satisfying
$b_{r^*}=\overline{b_r}$.

Given $\eps,\tau,\kappa>0$ and any neighborhood of the physical map
and its inverse in the smooth topology, there are a finite sequence of
passages, a final intermediate map $S'$, final packets $F'_r$, and a
smooth function $\Psi:\T\to\C$, such that
\begin{align}
 \sum_r\|F'_r-F_r\|_2&<\eps,
 \label{eq:prediction-transport}\\
 \int_\T\sqrt{p'(z)}\,\dd z
 &\ge \int_\T\sqrt{p(z)}\,\dd z-\tau,
 \qquad p'=\sum_r p'_r,
 \label{eq:prediction-hellinger}\\
 \left\|\sum_r b_rF'_r-
       \Psi(U_{S'})\sum_r F'_r\right\|_2&<\eps.
 \label{eq:prediction-error}
\end{align}
Here the differences in \eqref{eq:prediction-transport} are measured
in the original physical $L^2$ space.  All the physical maps and their
inverses can be kept in the prescribed neighborhood.  The final good
arc $I'$ satisfies $|\T\setminus I'|<\kappa$.
The output colors again have disjoint nonzero angular-index sets,
moduli independent of the spare coordinate, and conjugation symmetry.
The multiplier can be chosen with
$\Psi(-z)=\overline{\Psi(z)}$.
\end{proposition}

We first record the precise averaging assertion needed in the proof.
In the following lemma, a Laurent polynomial in $w_1,\ldots,w_n$ means
a finite linear combination of the monomials
$w^\nu=\prod_i w_i^{\nu_i}$, with $\nu\in\Z^n$.

\begin{lemma}[Finite phase tests]
\label{lem:phase-tests}
Fix finitely many expressions
\[
 P_j(z,w)=\sum_{\nu\in\mathcal E_j}c_{j,\nu}(z)w^\nu,
 \qquad c_{j,\nu}\in L^1(\T),
\]
where each $\mathcal E_j$ is finite.  For every $\eta>0$, successively
large positive integers $L_1,\ldots,L_n$ can be chosen so that,
simultaneously for all $j$,
\begin{equation}
 \left|\int_\T P_j\bigl(z,\ee(L_1z),\ldots,\ee(L_nz)\bigr)\,\dd z
 -\int_\T\int_{\T^n}P_j\bigl(z,\ee(y_1),\ldots,\ee(y_n)\bigr)
             \,\dd y\,\dd z\right|<\eta.
 \label{eq:phase-tests}
\end{equation}
At each selection, $L_i$ may additionally be required to exceed any
finite lower bound determined by the preceding choices and the
then-current construction data.
\end{lemma}

\begin{proof}
Let $D$ bound the absolute values of all the exponents in all the
$\mathcal E_j$.  The integral of a nonconstant monomial on the product
torus is zero.  On the one-dimensional diagonal its contribution is a
Fourier coefficient of $c_{j,\nu}$ at
$-\sum_i\nu_iL_i$.
The Fourier coefficients of an $L^1$ function tend to zero: this
follows by approximation in $L^1$ by trigonometric polynomials, whose
sufficiently high Fourier coefficients vanish.  Since the coefficient
list is finite, choose $R$ so that the sum of the bounds for all these
nonconstant contributions is less than $\eta$ whenever their
frequencies have absolute value greater than $R$.
Choose successively
\[
 L_j>R+D\sum_{i<j}L_i,
\]
as well as the additional finite lower bounds.  If $j$ is the last
nonzero position of $\nu$, then
\[
 \left|\sum_i\nu_iL_i\right|
 \ge L_j-D\sum_{i<j}L_i>R.
\]
This proves the assertion.  No independence assertion about the
distribution of the entire diagonal is made or needed.
\end{proof}

\begin{proof}[Proof of Proposition~\ref{prop:prediction}]
Zero colors may be omitted.  If all colors vanish, a neutral passage
with zero output packets and $\Psi=0$ proves the assertion, including
the requested new good arc.  Hence assume that a nonzero color is
present.
Dividing all tags by $\max(1,\max_r|b_r|)$ and reducing the requested
prediction tolerance reduces the proof to $|b_r|\le1$.  The transport
and square-root tolerances remain independently adjustable.  We work
under this normalization until the end.

\paragraph{The data fixed at the beginning of a block.}
For the moment suppose that the current packet mass is at most $2$.
Write its actual densities as $p_r$, put $p=\sum_rp_r$, and on
$\{p>0\}$ set
\[
 t_r=\frac{p_r}{p},\qquad
 \beta=\sum_r t_rb_r.
\]
Set $\beta=0$ on $\{p=0\}$.  Then $|\beta|\le1$ and
$\beta(-z)=\overline{\beta(z)}$.  Fix $0<\delta\le1/100$.
There is a smooth circle function $q$ such that
\begin{equation}
 |q|\le2,\qquad q(-z)=\overline{q(z)},\qquad
 \int_\T\sqrt p\,|\beta-q|^2\,\dd z<\delta.
 \label{eq:block-centering}
\end{equation}
Indeed, convolve $\beta$ with an even nonnegative smooth approximate
identity.  The convolutions have absolute value at most $1$, have
the required symmetry, and converge almost everywhere to $\beta$.
Since $\sqrt p\in L^1$, dominated convergence gives
\eqref{eq:block-centering}.  In particular this choice does not require
any uniform bound on derivatives of $q$.

Let $A$ be any fixed bounded function on the temporal circle.  In the
iteration, $A$ will be the sum of the estimators from $k$ already
completed blocks.  Both $q$ and $A$ are fixed before any frequencies
for the present block are chosen.  The densities $p_r$ in every model
below are these actual \emph{pre-block} densities.

\paragraph{Anchoring, with the number of passages fixed first.}
Choose an integer $c_0$ large enough that the full compact ranges
$m(c_0+\phi(\T))$ for distinct old indices are separated and avoid
zero, as required by Lemma~\ref{lem:passage}.  In particular this
separates the compact packet images
\[
 m(c_0+\phi(u)),\qquad
 u\in\operatorname{proj}_u\supp A_{r,m},
\]
for all distinct old indices $m$.  Let $K_r$ be the union of the
images belonging to color $r$.
The $K_r$ are disjoint compact sets and
$K_{r^*}=-K_r$.  Choose slightly larger compact neighborhoods of them,
still disjoint and off zero, and then disjoint open neighborhoods with
compact closures.  We will retain enough room between these two
families of neighborhoods for all the support movements in the block.
There is a smooth compactly supported function on $\R\setminus\{0\}$
such that
\begin{equation}
 d(-x)=\overline{d(x)},\qquad
 d(x)=b_r-q(x\bmod1)
 \quad\hbox{on a neighborhood of the enlarged }K_r.
 \label{eq:block-band-function}
\end{equation}
For example, use disjoint smooth cutoffs equal to one near these sets,
chosen compatibly with reflection, and multiply the $r$th cutoff by
$b_r-q(x\bmod1)$.  Fix $0<b_0<D_0$ such that all the retained
neighborhoods lie strictly inside $\{b_0<|x|<D_0\}$.

Apply Lemma~\ref{lem:signal} to this fixed $d$ and these band bounds.
It gives a modulus $P$ and an arbitrarily small positive strength
$\gamma$, independently of the future angular scale.  Put
\[
 \lambda=\gamma\frac{|(\Z/P\Z)^\times|}{P}.
\]
Take $\gamma$ so small that $6\lambda<1$ and $\lambda^2\le\delta/2$,
and choose an integer $n\ge1$ with
\begin{equation}
 \frac\delta2\le n\lambda^2\le\delta.
 \label{eq:block-budget}
\end{equation}
One may take $n=\lfloor\delta/\lambda^2\rfloor$.
At this point $q,c_0,K_r,d,P,\gamma,\lambda,n$ are all fixed.
In particular they have been chosen \emph{before} the integer of the
anchor passage.  The error targets used below depend only on these
data, the already fixed $A,k$, and the requested tolerance $\nu$.
They can therefore be fixed now, before executing the anchor.

Now make a neutral passage with the already fixed $c_0$, choosing its
jitter and its integer $N_0$ for the required accuracy.  Use
$B=N_0$ as the output angular scale.  Lemmas~\ref{lem:passage}
and~\ref{lem:neutral} give color densities as close to $p_r$ in
$L^1$ as desired.  On their retained supports the normalized indices
$x=j/N_0$ lie in arbitrarily small enlargements of $K_r$.  If $J_0$
is the bounded jitter chosen for this passage, their actual temporal
frequencies satisfy
\[
 z=x+xJ_0(a)/N_0\pmod1.
\]
Thus $\dist_{\T}(x\bmod1,z)$ can be made arbitrarily small after
$J_0$ is fixed, by increasing $N_0$.

We follow this anchor by $n$ refresh--transmit pairs from
Lemma~\ref{lem:refresh-transmit}. At every passage the scale changes by
$B'=NcB$. The support estimates
\eqref{eq:signal-index-drift}--\eqref{eq:signal-temporal-drift}, which
apply to neutral as well as transmitting passages, control normalized-index
and temporal-frequency movement separately. Choose $c$ and sufficiently
small support enlargements to meet the movement budgets; in particular,
the enlargement is smaller than the temporal tolerance because its
contribution there is not divided by $N$. Fix the bounded jitter $J$
next, and choose $N$ last for the remaining terms. Thus an arbitrarily
broad neutral jitter remains compatible with temporal tracking.

There are only $2n+1$ passages in this block.  Apportion small enough
movement budgets among them so that all normalized indices remain
in the neighborhoods on which \eqref{eq:block-band-function} holds,
and so that $\dist_{\T}(x\bmod1,z)$ remains below any prescribed
positive tolerance.  Uniform continuity of the fixed $q$ then makes
\begin{equation}
 |d(x)-(b_r-q(z))|
 \label{eq:block-substitution}
\end{equation}
uniformly as small as needed on every retained support of color $r$.
The positive gaps preserve disjoint angular-index supports.  Each
passage preserves the independent-of-spare-coordinate modulus
property and conjugation symmetry.

\paragraph{The product model and its approximation.}
Put $a_r(z)=b_r-q(z)$, so $|a_r|\le3$.  Denote the integer of the
$i$th transmission by $L_i$.  The model after $i$ transmissions is
\begin{equation}
 h_r^{(i)}(z)=p_r(z)\prod_{j=1}^i f_{r,j}(z),\qquad
 f_{r,j}(z)=1+2\lambda\Rea\{a_r(z)\ee(-L_jz)\}.
 \label{eq:block-model}
\end{equation}
All factors are positive, and $f_{r,j}\le M=1+6\lambda<2$.
Let $g_r^{(i)}$ be the actual densities after these pairs, with
$g_r^{(0)}$ the anchor densities.  The refresh--transmit lemma,
including the substitution \eqref{eq:block-substitution}, allows
\[
 \sum_r\|g_r^{(i)}-g_r^{(i-1)}f_{r,i}\|_1<\eta_i,
 \qquad
 \sum_r\|g_r^{(0)}-p_r\|_1<\eta_0.
\]
Indeed the substitution costs at most
$2\gamma\omega\sum_r\int g_r^{(i-1)}$ if its uniform error is
$\omega$; the remaining error is freely adjustable.  We may first
choose the movement tolerance small enough for all the finitely many
$\eta_i$, and then implement each pair on its actual finite input.
Consequently
\begin{equation}
 E_i:=\sum_r\|g_r^{(i)}-h_r^{(i)}\|_1
 \le ME_{i-1}+\eta_i,
 \qquad
 E_n\le M^n\eta_0+\sum_{i=1}^nM^{n-i}\eta_i.
 \label{eq:block-error-propagation}
\end{equation}
For any prescribed $e_0>0$, it suffices to choose every
$\eta_i<e_0/((n+1)M^n)$ to obtain $E_n<e_0$.
All quantities in this accuracy requirement were fixed before the
anchor.  Even if $\lambda$ is extremely small and $n$ and $M^n$ are
extremely large, they are finite known numbers.

The frequencies $L_i$ remain free to exceed arbitrary successive
lower bounds.  In particular, Lemma~\ref{lem:phase-tests} can be
applied to any finite list of Laurent polynomials whose coefficient
functions depend only on the pre-block data
\begin{equation}
 p_r,\ t_r,\ q,\ A,\ b_r,\ n,\ \lambda.
 \label{eq:block-fixed-data}
\end{equation}
The actual packets depend on previous choices, but appear only in
\eqref{eq:block-error-propagation}, not in this fixed coefficient
list.  At the $i$th pair, first choose its refresh and all its
transmission data preceding the passage integer; then choose $L_i$
large enough for both the physical construction and the phase-test
lower bound.  Thus the averaging and the construction impose
compatible lower bounds on the same integer.

\paragraph{The square-root integral.}
Introduce independent Haar phases $Z_1,\ldots,Z_n$ solely for
computing the product-torus averages in
Lemma~\ref{lem:phase-tests}.  At each fixed $z$ put
\[
 H_r=\prod_{i=1}^n
       [1+\lambda(a_r\overline{Z_i}+\overline{a_r}Z_i)],
 \qquad R=\sum_r t_rH_r.
\]
The following pointwise calculations are on $\{p>0\}$; all their
integrands are multiplied by $p_r$ or $\sqrt p$ and are extended by
zero on its complement.  We have
\begin{equation}
 \mathbb E H_r=1,
 \qquad
 \mathbb E(H_rH_s)
       =(1+2\lambda^2\Rea(a_r\overline{a_s}))^n,
 \qquad \mathbb E R=1.
 \label{eq:block-second-moment-identity}
\end{equation}
For real $y$ and integer $n\ge1$, the binomial expansion gives
\[
 |(1+y)^n-1-ny|
 \le\sum_{j=2}^n\frac{(n|y|)^j}{j!}
 \le\frac{(n|y|)^2}{2}\exp(n|y|).
\]
Apply this with $y=2\lambda^2\Rea(a_r\overline{a_s})$;
then $n|y|\le18\delta$.  Since $\sum_rt_r=1$, the sum of the
linear terms in \eqref{eq:block-second-moment-identity} is
$2n\lambda^2|\sum_rt_ra_r|^2$.  Therefore
\begin{equation}
 \mathbb E(R-1)^2
 \le2\delta|\beta-q|^2+162e^{18\delta}\delta^2.
 \label{eq:block-second-moment-bound}
\end{equation}
This bound is uniform in the number of colors and in all the later
geometric parameters.

For every $x\ge0$,
\begin{equation}
 \sqrt x\ge1+\frac{x-1}{2}-(x-1)^2.
 \label{eq:block-square-root-bound}
\end{equation}
In fact, on setting $t=\sqrt x$, the difference between the two sides
equals $(t-1)^2(t^2+2t+1/2)$.  Multiplying
\eqref{eq:block-square-root-bound} by $\sqrt p$, taking the
independent-phase expectation, and using
\eqref{eq:block-centering} and \eqref{eq:block-second-moment-bound}
shows that the averaged polynomial lower bound is at least
\[
 \int_\T\sqrt p\,\dd z
 -2\delta^2-162e^{18\delta}\delta^2\int_\T\sqrt p\,\dd z
 \ge \int_\T\sqrt p\,\dd z-300\delta^2.
\]
Here $\int\sqrt p\le\sqrt2$ and $\delta\le1/100$ were used.

The expression
\[
 \sqrt p\,[1+(R-1)/2-(R-1)^2]
\]
is a Laurent polynomial in the phases with fixed integrable
coefficients.  Lemma~\ref{lem:phase-tests} therefore transfers this
lower bound, up to any prescribed error, to the diagonal phases
$Z_i=\ee(L_i z)$.  On that diagonal
$pR=\sum_r h_r^{(n)}$.
There is no need to apply the lemma to $\sqrt R$ itself.
Finally, for nonnegative densities $g,h$ on the unit circle,
\begin{equation}
 \left|\int_\T\sqrt g-\int_\T\sqrt h\right|
 \le\left(\int_\T|g-h|\right)^{1/2},
 \label{eq:block-density-hellinger-error}
\end{equation}
by $(\sqrt g-\sqrt h)^2\le|g-h|$ and Cauchy--Schwarz.
Thus the actual output density $p^{\rm out}=\sum_rg_r^{(n)}$
satisfies, for any chosen $\nu>0$ and sufficiently accurate
construction,
\begin{equation}
 \int_\T\sqrt{p^{\rm out}}\,\dd z
 \ge\int_\T\sqrt p\,\dd z-300\delta^2-\nu.
 \label{eq:block-hellinger-step}
\end{equation}

\paragraph{Prediction and a previous estimator.}
With the same formal phases, put
\[
 D=q+\frac1{n\lambda}\sum_{i=1}^n Z_i.
\]
For a fixed color $r$ and fixed $z$, $H_r$ is a positive probability
density on the product phase torus, and under this density the phases
remain independent.  Direct integration of one factor gives
\[
 \mathbb E_{H_r}Z_i=\lambda a_r,
 \qquad
 \mathbb E_{H_r}|Z_i-\lambda a_r|^2
       =1-\lambda^2|a_r|^2.
\]
Consequently
\begin{equation}
 \mathbb E_{H_r}D=b_r,
 \qquad
 \mathbb E_{H_r}|D-b_r|^2
 =\frac{1-\lambda^2|a_r|^2}{n\lambda^2}
 \le\frac2\delta.
 \label{eq:block-estimator-moments}
\end{equation}
For the arbitrary fixed bounded previous estimator $A$, this means
that the cross term against $A-kb_r$ vanishes \emph{at each $z$}:
\[
 \mathbb E_{H_r}|A+D-(k+1)b_r|^2
 =|A-kb_r|^2+\mathbb E_{H_r}|D-b_r|^2.
\]
After multiplying by $p_r$, summing, and integrating, the increase
in this mean square is at most $4/\delta$.

The expression
\[
 \sum_r p_rH_r|A+q+(n\lambda)^{-1}\textstyle\sum_iZ_i-(k+1)b_r|^2
\]
is another finite Laurent polynomial with the fixed integrable
coefficients in \eqref{eq:block-fixed-data}.  Apply
Lemma~\ref{lem:phase-tests} to it simultaneously with the preceding
square-root lower-bound test.  Write
\begin{equation}
 D(z)=q(z)+\frac1{n\lambda}\sum_{i=1}^n\ee(L_i z)
 \label{eq:block-diagonal-estimator}
\end{equation}
for the resulting actual multiplier.  Its norm is bounded by
$2+1/\lambda$, independently of the choices of $L_i$.  In passing
from model to actual color densities, the error in the prediction
integral is at most $W E_n$, where
\begin{equation}
 W=(\|A\|_\infty+2+1/\lambda+k+1)^2.
 \label{eq:block-estimator-error-weight}
\end{equation}
This finite number is known before choosing the anchor accuracy or
any transmission frequencies.  For example, choose
\[
 e_0<\min\{(\nu/4)^2,\nu/(4W)\},
\]
and make the two phase-test errors less than $\nu/2$.
Together with \eqref{eq:block-error-propagation} this proves
\begin{equation}
 \sum_r\int_\T p_r^{\rm out}|A+D-(k+1)b_r|^2\,\dd z
 \le\sum_r\int_\T p_r|A-kb_r|^2\,\dd z
            +\frac4\delta+\nu,
 \label{eq:block-prediction-step}
\end{equation}
simultaneously with \eqref{eq:block-hellinger-step}.
The same $e_0$ and support accuracies can be imposed before the
anchor.  Thus a small strength entails a more demanding finite
accuracy requirement, without a circular parameter choice.

\paragraph{Iteration of blocks.}
All color vectors can be transported with an arbitrarily small sum
of physical norm errors at every passage.  Reserve a total such
budget smaller than
\[
 \min\{\eps,\sqrt2-\sqrt{3/2}\}.
\]
Since the sum of the color errors bounds the error of their sum,
all intermediate packet sums then have squared norm less than $2$.
This supplies the mass bound used above at every step, with room to
spare.  Smooth map and inverse errors can likewise be apportioned
among the finitely many blocks and, after $n$ is chosen in a block,
among its $2n+1$ passages.  Choosing a basic neighborhood specified
by finitely many smooth norms inside the requested neighborhood
makes this allocation explicit.  Every passage may also make the
new excluded arc arbitrarily small; apply this in particular at the
last passage.

Set $\nu\le\delta^2$ in each block.  Start with $A_0=0$, and let
$A_{k+1}=A_k+D_{k+1}$ after block $k+1$.  At the beginning of each
new block choose a \emph{new} $q$ from that block's actual input
densities.  That $q$ is then held fixed throughout its block.
The band function, sieve modulus, strength, and length may all
change from block to block.  The estimates depend only on
$|b_r|\le1$, $|q|\le2$, the mass bound $2$, and
\eqref{eq:block-budget}; they require no uniform lower bound for
the strengths and no uniform derivative bound for the successive
centering functions.

After $K$ blocks, telescoping
\eqref{eq:block-hellinger-step} and
\eqref{eq:block-prediction-step} gives
\begin{align}
 \int_\T\sqrt{p^{(K)}}\,\dd z
 &\ge\int_\T\sqrt{p^{(0)}}\,\dd z-400K\delta^2,
 \label{eq:block-total-hellinger}\\
 \sum_r\int_\T p_r^{(K)}
       |A_K/K-b_r|^2\,\dd z
 &\le\frac5{K\delta}.
 \label{eq:block-total-prediction}
\end{align}
For the second estimate the initial error is zero, and
$4/\delta+\delta^2\le5/\delta$.
Given the desired prediction tolerance $\eps$, first choose
$\delta>0$ sufficiently small and then put
\[
 K=\left\lceil\frac{10}{\eps^2\delta}\right\rceil.
\]
We can require simultaneously $\delta\le1/100$ and
\[
 400\left(\frac{10\delta}{\eps^2}+\delta^2\right)<\tau.
\]
Then $K\delta$ is large enough for
\eqref{eq:block-total-prediction}, and $K\delta^2$ is small enough
for \eqref{eq:block-total-hellinger}.  These choices precede all
blocks and do not depend on their future strengths or lengths.

Set $\Psi=A_K/K$.  Disjoint angular-index supports make all
different color cyclic subspaces orthogonal, so the spectral
density formula gives exactly
\[
 \left\|\sum_r b_rF_r^{(K)}-
          \Psi(U_{S^{(K)}})\sum_rF_r^{(K)}\right\|_2^2
 =\sum_r\int_\T p_r^{(K)}|b_r-\Psi|^2\,\dd z.
\]
This proves the required norm estimate.  Every $D$ in
\eqref{eq:block-diagonal-estimator} is smooth and has conjugation
symmetry, hence so does $\Psi$.  Restoring the original scale of
the tags proves the proposition in its stated form.
\end{proof}
